% ============================================================
% BÖLÜM 1: GİRİŞ
% Genişletilmiş ve Akademik Referanslarla Zenginleştirilmiş
% ============================================================

\chapter{GİRİŞ}

\section{Araştırmanın Bağlamı}

Psikolojik danışmanlık ve rehberlik (PDR), bireylerin psikolojik sağlığını korumak, geliştirmek ve sorunlarını çözmelerine yardımcı olmak amacıyla yürütülen profesyonel bir yardım sürecidir. Bu süreç, özellikle çocukluk döneminde kritik bir öneme sahiptir; zira erken dönemde tespit edilen psikolojik sorunlar, uygun müdahalelerle daha etkili biçimde ele alınabilmektedir (Kuzgun, 2000; Yeşilyaprak, 2016). Okul psikolojik danışmanları ve rehber öğretmenler, bu bağlamda çocukların gelişimsel ihtiyaçlarını karşılamak, potansiyel sorunları erken dönemde tespit etmek ve gerekli yönlendirmeleri yapmakla sorumludur.

\subsection{Çocuk Psikolojisi ve Değerlendirme Zorlukları}

Çocuklarla yürütülen psikolojik değerlendirme süreçlerinde, yetişkinlere uygulanan standart değerlendirme araçlarının kullanımı çoğu zaman sınırlı kalmaktadır. Piaget'nin (1952) bilişsel gelişim kuramına göre, çocuklar soyut düşünme becerilerini ancak 11-12 yaşlarından itibaren geliştirmeye başlamaktadır. Bu durum, sözel ifadeye dayalı değerlendirme araçlarının etkinliğini sınırlamaktadır.

Çocuklar, duygularını ve düşüncelerini sözel olarak ifade etmede zorluk yaşayabilmekte; bu durum, psikolojik değerlendirme süreçlerini karmaşıklaştırmaktadır (Malchiodi, 1998). Erikson'un (1963) psikososyal gelişim kuramına göre, çocukluk döneminde yaşanan krizlerin çözümü, gelecekteki psikolojik sağlık için belirleyici öneme sahiptir. Bu nedenle, çocukların iç dünyalarını anlamak için alternatif değerlendirme yöntemlerine ihtiyaç duyulmaktadır.

Bu noktada, projektif testler ve özellikle çizim temelli değerlendirme araçları, çocukların iç dünyalarına ulaşmada önemli bir pencere sunmaktadır. Winnicott (1971), çizimlerin çocuklar için güvenli bir geçiş nesnesi işlevi gördüğünü ve içsel deneyimlerin dışavurumu için doğal bir araç oluşturduğunu belirtmiştir.

\subsection{Projektif Çizim Testlerinin Önemi}

House-Tree-Person (HTP) testi, John Buck tarafından 1948 yılında geliştirilmiş ve çocukların ev, ağaç ve insan figürü çizmelerini içeren projektif bir değerlendirme aracıdır (Buck, 1948). Bu test, çocukların bilinçaltı düşüncelerini, duygularını, aile ilişkilerini ve benlik algılarını yansıttığı varsayılan sembolik ifadeler içermektedir. HTP testi, dünya genelinde yaygın biçimde kullanılmakta ve özellikle çocuk değerlendirmesinde önemli bir yer tutmaktadır (Hammer, 1958; Jolles, 1971).

Emanuel Hammer (1958), HTP testinin klinik uygulamasını genişletmiş ve çizimlerdeki sembolik anlamların sistematik yorumlanması için kapsamlı bir çerçeve sunmuştur. Karen Machover (1949) ise insan figürü çizimlerinin kişilik değerlendirmesindeki rolünü vurgulayan ``Draw-A-Person'' (DAP) testini geliştirmiştir. Charles Koch (1952), ağaç çizimlerinin bilinçdışı benlik algısını yansıttığını ileri sürerek ``Ağaç Testi'' yaklaşımını ortaya koymuştur.

\subsection{Teknolojinin Psikolojik Danışmanlıkta Kullanımı}

Ancak HTP testinin yorumlanması, uzun yıllar süren eğitim ve deneyim gerektiren karmaşık bir süreçtir. Groth-Marnat (2009), projektif testlerin yorumlanmasının klinisyenin teorik altyapısına, deneyimine ve kültürel farkındalığına bağlı olduğunu vurgulamıştır. Okul ortamlarında görev yapan psikolojik danışmanların çoğu, yoğun iş yükleri ve sınırlı zaman kaynakları nedeniyle bu tür kapsamlı değerlendirmeleri gerçekleştirmekte güçlük çekmektedir. Ayrıca, projektif testlerin yorumlanmasında öznel yargıların etkisi, değerlendirmeler arası tutarlılık sorunlarına yol açabilmektedir (Lilienfeld, Wood ve Garb, 2000).

Bu bağlamda, teknolojinin psikolojik danışmanlık hizmetlerine entegrasyonu, hem verimliliği artırma hem de değerlendirme süreçlerini standartlaştırma potansiyeli taşımaktadır. Yapay zekâ teknolojilerindeki son gelişmeler, özellikle çok modlu büyük dil modelleri (Multimodal Large Language Models - MLLM), görsel içeriklerin analizinde ve doğal dil işlemede önemli ilerlemeler kaydetmiştir (Brown ve ark., 2020; OpenAI, 2023). Bu teknolojiler, çocuk çizimlerinin otomatik analizinde yeni olanaklar sunmaktadır.

\section{Problemin Tanımı}

Okul psikolojik danışmanlığı alanında çocuk çizimlerinin değerlendirilmesi sürecinde çeşitli sorunlar yaşanmaktadır:

\begin{enumerate}
  \item \textbf{Uzmanlık Gereksinimleri:} HTP testinin doğru yorumlanması, projektif testler konusunda kapsamlı eğitim ve deneyim gerektirmektedir. Okul psikolojik danışmanlarının çoğu, lisans eğitimleri sırasında bu konuda yeterli derinlikte eğitim almamaktadır.

  \item \textbf{Zaman Kısıtlamaları:} Okul ortamlarında danışman başına düşen öğrenci sayısının yüksekliği, bireysel değerlendirmelere ayrılabilecek zamanı kısıtlamaktadır. Türkiye'de bir okul psikolojik danışmanına düşen öğrenci sayısı, uluslararası standartların çok üzerindedir (MEB, 2020).

  \item \textbf{Değerlendirmeler Arası Tutarsızlık:} Projektif testlerin öznel yorumlama içermesi, farklı değerlendiricilerin aynı çizim için farklı sonuçlara ulaşmasına neden olabilmektedir (Anastasi ve Urbina, 1997).

  \item \textbf{Güncel Literatüre Erişim:} Akademik literatürdeki gelişmelerin takip edilmesi ve yorumlama süreçlerine entegre edilmesi zaman ve kaynak gerektirmektedir.

  \item \textbf{Erken Tespit Zorlukları:} Psikolojik sorunların erken dönemde tespit edilmesi, uygun müdahalelerin zamanında yapılabilmesi için kritiktir. Ancak sınırlı kaynaklar, kapsamlı tarama çalışmalarını zorlaştırmaktadır (Costello ve ark., 2003).
\end{enumerate}

\section{Çalışmanın Amacı}

Bu çalışmanın temel amacı, okul psikolojik danışmanlarının çocuk çizimlerini değerlendirme süreçlerini destekleyecek, yapay zekâ tabanlı bir karar destek sistemi geliştirmektir. NIDA (Neural Imagery Decision Assessment) adı verilen bu sistem, aşağıdaki hedefleri gerçekleştirmeyi amaçlamaktadır:

\begin{enumerate}
  \item HTP çizimlerinin akademik literatüre dayalı, tutarlı ve kaynak destekli yorumlarını sunmak
  \item Psikolojik danışmanların değerlendirme süreçlerindeki iş yükünü azaltmak
  \item Erken dönem psikolojik sorunların tespitinde yardımcı bir araç sağlamak
  \item PDR alanında teknoloji entegrasyonuna örnek bir uygulama oluşturmak
  \item Akademik literatürdeki yorumlama kurallarını yapılandırılmış ve erişilebilir bir formata dönüştürmek
\end{enumerate}

\section{Çalışmanın Kapsamı}

Bu çalışma, aşağıdaki kapsam dahilinde yürütülmüştür:

\begin{itemize}
  \item \textbf{Yaş Grubu:} 5-12 yaş arası çocuklar (ilkokul dönemi). Bu yaş aralığı, Lowenfeld'in (1947) çizim gelişim evrelerinden Ön-Şematik, Şematik ve Gerçekçilik evrelerini kapsamaktadır.
  \item \textbf{Çizim Türü:} House-Tree-Person (HTP) testi - Ev, Ağaç, İnsan ve Tam HTP çizimleri
  \item \textbf{Hedef Kullanıcılar:} Okul psikolojik danışmanları ve rehber öğretmenler
  \item \textbf{Teknoloji:} Web tabanlı uygulama, Gemini 2.0 Flash çok modlu dil modeli
  \item \textbf{Dil Desteği:} Türkçe (varsayılan) ve İngilizce
\end{itemize}

\section{Araştırmanın Sınırlılıkları}

Bu çalışma, aşağıdaki sınırlılıklar çerçevesinde değerlendirilmelidir:

\begin{enumerate}
  \item \textbf{Tanı Koyma Kapasitesi:} NIDA, kesinlikle bir tanı aracı değildir ve profesyonel değerlendirmenin yerini almaz. Sistem, yalnızca karar destek işlevi görmektedir. Amerikan Psikoloji Derneği'nin (APA, 2017) etik kuralları, tanı koyma yetkisinin yalnızca lisanslı profesyonellere ait olduğunu vurgulamaktadır.

  \item \textbf{Akademik Geçerlilik:} Projektif çizim testlerinin geçerliliği akademik literatürde tartışmalıdır (Lilienfeld ve ark., 2000; Hunsley ve Bailey, 1999). NIDA, bu tartışmayı çözmemekte, mevcut literatürü şeffaf biçimde sunmaktadır.

  \item \textbf{Kültürel Bağlam:} HTP yorumlama kurallarının çoğu Batı kültürlerinde geliştirilmiştir. Türk kültürel bağlamına uyarlama çalışmaları sınırlıdır (Yavuzer, 2009). Kültürlerarası geçerlilik çalışmaları gelecekte yapılmalıdır.

  \item \textbf{Yapay Zekâ Sınırlılıkları:} Büyük dil modelleri halüsinasyon üretebilir ve tutarsız sonuçlar verebilir (Bender ve ark., 2021). Sistem, bu riskleri azaltmak için bilgi tabanı sınırlaması uygulamaktadır.

  \item \textbf{Veri Seti Eksikliği:} Çalışma, özel bir veri seti ile eğitilmiş bir model içermemekte; tamamen akademik literatüre dayalı kural tabanlı bir yaklaşım benimsemektedir.
\end{enumerate}

\newpage
